%=====================================================================
\chapter{Machine Learning Fundamentals}\label{ch:mlfundamentals}
%=====================================================================
\locator{3}{Part III --- Machine Learning Core}

\begin{outcomes}
By the end of this chapter you will be able to:
\begin{tight}
  \item \textbf{Distinguish} supervised, unsupervised, semi-supervised and reinforcement learning, and match a problem to the right one.
  \item \textbf{Explain} the purpose of the train/validation/test split and why the test set may be used only once.
  \item \textbf{Diagnose} overfitting and underfitting from a learning curve.
  \item \textbf{Describe} the bias--variance trade-off and how model complexity moves along it.
  \item \textbf{Justify} why a baseline must precede any model.
\end{tight}
\end{outcomes}

\begin{prereq}
\chref{math} (vectors, loss, gradient descent), \chref{wrangling} (splitting and
leakage) and \chref{descriptive}.
\end{prereq}

\begin{keyterms}
supervised / unsupervised / reinforcement learning $\bullet$ label $\bullet$
training, validation and test sets $\bullet$ generalisation $\bullet$ overfitting
$\bullet$ underfitting $\bullet$ bias $\bullet$ variance $\bullet$ learning curve
$\bullet$ hyperparameter $\bullet$ baseline $\bullet$ no free lunch
\end{keyterms}

\section{What Learning Means Here}

\begin{definitionbox}[Machine Learning, Operationally]
A program learns from experience $E$ with respect to task $T$ and performance
measure $P$ if its performance on $T$, measured by $P$, improves with $E$.
(Mitchell, 1997.) In practice: $E$ is your training data, $T$ is the prediction
you framed in \chref{lifecycle}, and $P$ is the metric of \chref{evaluation}.
\end{definitionbox}

The crucial word is \term{generalisation}. Memorising the training data is trivial
and worthless; the goal is performance on data the model has never seen. Everything
in this chapter exists to protect that distinction.

\section{The Four Paradigms}

\dsfig{%
\begin{tikzpicture}[
  box/.style={rounded corners=4pt, draw=#1, line width=1.1pt, fill=#1!8,
              text width=3.15cm, align=center, font=\scriptsize, inner sep=5pt},
  child/.style={rounded corners=2pt, draw=#1, fill=#1!5, text=#1!50!black,
                text width=2.9cm, align=center, font=\tiny, inner sep=3pt},
  arr/.style={-{Stealth[length=2.2mm]}, primary!50, line width=0.9pt}]

\node[box=primary] (root) at (5.85,1.5) {{\color{primary}\bfseries\small \faIcon{brain}~Machine Learning}\\[2pt]
       {\tiny how much supervision is available?}};

\node[box=secondary] (sup) at (0,-0.9) {{\color{secondary}\bfseries SUPERVISED}\\[2pt]
      labelled data\\$(x_i, y_i)$};
\node[box=goldhl] (semi) at (3.9,-0.9) {{\color{goldhl!60!black}\bfseries SEMI-SUPERVISED}\\[2pt]
      a few labels,\\many unlabelled};
\node[box=greenhl] (uns) at (7.8,-0.9) {{\color{greenhl!50!black}\bfseries UNSUPERVISED}\\[2pt]
      no labels at all};
\node[box=purplehl] (rl) at (11.7,-0.9) {{\color{purplehl}\bfseries REINFORCEMENT}\\[2pt]
      rewards from\\an environment};

\foreach \n in {sup,semi,uns,rl}{\draw[arr] (root) -- (\n);}

\node[child=secondary] (cls) at (-1.6,-2.9) {Classification\\\emph{will this fail?}};
\node[child=secondary] (reg) at (1.6,-2.9)  {Regression\\\emph{how long until?}};
\node[child=greenhl] (clu) at (6.2,-2.9)    {Clustering\\\emph{which groups?}};
\node[child=greenhl] (dim) at (9.4,-2.9)    {Dimensionality\\reduction, anomalies};
\draw[arr] (sup) -- (cls); \draw[arr] (sup) -- (reg);
\draw[arr] (uns) -- (clu); \draw[arr] (uns) -- (dim);

\node[font=\tiny\itshape, text=gray!55!black, text width=3.2cm, align=center] at (3.9,-2.9)
     {labelling is the expensive part; this is why it matters};
\node[font=\tiny\itshape, text=gray!55!black, text width=3.0cm, align=center] at (11.7,-2.9)
     {covered in \chref{agentic}, not Part III};
\end{tikzpicture}}%
{The learning paradigms, organised by how much supervision the data provides. Chapters
12--15 cover the supervised branch; \chref{unsupervised} covers the unsupervised
one.}{ml-paradigms}

\begin{conceptbox}[The Question That Picks Your Paradigm]
\emph{Do I have examples of the right answer?} If yes, supervised. If no, but you
want structure, unsupervised. If you have a few answers and lots of unlabelled
data, semi-supervised. If the answer only emerges from acting and seeing what
happens, reinforcement.

In practice the binding constraint is almost never the algorithm --- it is whether
anyone will pay to label 5{,}000 examples.
\end{conceptbox}

\section{The Three-Way Split}

\begin{definitionbox}[Training, Validation and Test Sets]
\term{Training set} (typically 60--70\%): the model learns its parameters here.\\
\term{Validation set} (15--20\%): you compare models and tune hyperparameters here.\\
\term{Test set} (15--20\%): used \textbf{once}, at the very end, to estimate real-world
performance.
\end{definitionbox}

\dsfig{%
\begin{tikzpicture}[font=\scriptsize]
\node[rounded corners=3pt, fill=primary!16, draw=primary, minimum width=7.3cm,
      minimum height=0.75cm, text=primary!55!black, font=\scriptsize\bfseries] at (3.65,0) {TRAINING (65\%)};
\node[rounded corners=3pt, fill=goldhl!22, draw=goldhl, minimum width=2.2cm,
      minimum height=0.75cm, text=goldhl!45!black, font=\scriptsize\bfseries] at (8.4,0) {VALID.\ (17\%)};
\node[rounded corners=3pt, fill=accent!15, draw=accent, minimum width=2.2cm,
      minimum height=0.75cm, text=accent!70!black, font=\scriptsize\bfseries] at (10.9,0) {TEST (18\%)};

\node[align=center, text width=7.3cm, font=\tiny, text=gray!55!black] at (3.65,-1.0)
     {fit parameters here\\\emph{used thousands of times}};
\node[align=center, text width=2.5cm, font=\tiny, text=gray!55!black] at (8.4,-1.0)
     {choose model \& hyperparameters\\\emph{used dozens of times}};
\node[align=center, text width=2.5cm, font=\tiny, text=accent!70!black] at (10.9,-1.0)
     {\textbf{used ONCE}\\\emph{ever}};

\node[align=left, text width=13.0cm, font=\tiny, text=accent!70!black] at (4.9,-2.2)
     {\textbf{Why only once?} Every time you look at the test score and change something,
      you are fitting to it --- slowly, by hand. After twenty such peeks the test set has
      become a second validation set and no longer estimates anything honest. This is
      leakage through the researcher rather than through the code, and it is just as
      damaging.};
\draw[decorate, decoration={brace, amplitude=4pt, mirror}, accent!60]
     (9.8,-0.45) -- (12.0,-0.45);
\end{tikzpicture}}%
{The three-way split and the discipline it demands. The validation set exists precisely so
that the test set can stay untouched --- without it, every tuning decision contaminates
your final estimate.}{three-way-split}

\begin{alertbox}[Stratify, and Split by Entity]
When classes are imbalanced, use a \emph{stratified} split so each part keeps the
same class proportions --- otherwise a rare class may be absent from the test set
entirely. And when rows are grouped (many readings per device, many sessions per
user), split by \emph{device} or \emph{user}, not by row: the same entity appearing
in both train and test is leakage (\chref{wrangling}).
\end{alertbox}

\section{Overfitting, Underfitting and the Learning Curve}

\begin{definitionbox}[Overfitting and Underfitting]
\term{Overfitting}: the model has learned the noise as well as the signal. Training
error is low, validation error is high. \\
\term{Underfitting}: the model is too simple to capture the signal. Both errors are
high.
\end{definitionbox}

\dsfig{%
\begin{tikzpicture}
\begin{groupplot}[
  group style={group size=3 by 1, horizontal sep=0.95cm},
  width=5.35cm, height=4.2cm,
  axis lines=left, tick label style={font=\tiny}, label style={font=\tiny},
  title style={font=\scriptsize\bfseries},
  xmin=0, xmax=10, ymin=0, ymax=10, xtick=\empty, ytick=\empty,
]
\nextgroupplot[title={\color{secondary}Underfit (high bias)}]
\addplot[only marks, mark=*, mark size=1pt, gray!70] coordinates {
 (0.7,2.2)(1.4,4.1)(2.1,5.6)(2.8,6.9)(3.5,7.7)(4.2,8.2)(4.9,8.3)(5.6,8.0)
 (6.3,7.2)(7.0,6.1)(7.7,4.6)(8.4,3.0)(9.1,1.6)};
\addplot[secondary, line width=1.3pt, domain=0.4:9.4] {5.4};
\node[font=\tiny, secondary!60!black, align=center] at (axis cs:5,1.2) {misses the shape};

\nextgroupplot[title={\color{greenhl}Good fit}]
\addplot[only marks, mark=*, mark size=1pt, gray!70] coordinates {
 (0.7,2.2)(1.4,4.1)(2.1,5.6)(2.8,6.9)(3.5,7.7)(4.2,8.2)(4.9,8.3)(5.6,8.0)
 (6.3,7.2)(7.0,6.1)(7.7,4.6)(8.4,3.0)(9.1,1.6)};
\addplot[greenhl, line width=1.3pt, domain=0.4:9.4, samples=80]
  {-0.33*(x-4.75)^2 + 8.4};
\node[font=\tiny, greenhl!55!black, align=center] at (axis cs:5,1.2) {signal, not noise};

\nextgroupplot[title={\color{accent}Overfit (high variance)}]
\addplot[only marks, mark=*, mark size=1pt, gray!70] coordinates {
 (0.7,2.2)(1.4,4.1)(2.1,5.6)(2.8,6.9)(3.5,7.7)(4.2,8.2)(4.9,8.3)(5.6,8.0)
 (6.3,7.2)(7.0,6.1)(7.7,4.6)(8.4,3.0)(9.1,1.6)};
\addplot[accent, line width=1.1pt, smooth] coordinates {
 (0.7,2.2)(1.4,4.1)(2.1,5.6)(2.8,6.9)(3.5,7.7)(4.2,8.2)(4.9,8.3)(5.6,8.0)
 (6.3,7.2)(7.0,6.1)(7.7,4.6)(8.4,3.0)(9.1,1.6)};
\node[font=\tiny, accent!80!black, align=center] at (axis cs:5,1.2) {learns every wobble};
\end{groupplot}
\end{tikzpicture}}%
{The same 13 points fitted three ways. The overfitted curve passes through every point and
has zero training error --- and would predict badly for any new observation.}{over-underfit}

\dsfig{%
\begin{tikzpicture}
\begin{groupplot}[
  group style={group size=2 by 1, horizontal sep=1.5cm},
  width=6.6cm, height=4.7cm,
  axis lines=left, tick label style={font=\tiny}, label style={font=\tiny},
  title style={font=\scriptsize\bfseries},
  legend style={font=\tiny, draw=gray!40},
]
% --- error vs complexity
\nextgroupplot[title={Error against model complexity},
  xlabel={model complexity $\longrightarrow$}, ylabel={error},
  xmin=0, xmax=10, ymin=0, ymax=10, xtick=\empty, ytick=\empty,
  legend style={at={(0.5,0.99)}, anchor=north}]
\addplot[primary, line width=1.3pt, domain=0.4:9.6, samples=80] {7.5*exp(-0.42*x)+0.35};
\addlegendentry{training error}
\addplot[accent, line width=1.3pt, domain=0.4:9.6, samples=80]
  {7.5*exp(-0.42*x)+0.35 + 0.115*x^2};
\addlegendentry{validation error}
\draw[greenhl, dashed, line width=1pt] (axis cs:4.1,0) -- (axis cs:4.1,9.0);
\node[font=\tiny, greenhl!55!black, anchor=south] at (axis cs:4.1,9.0) {sweet spot};
\node[font=\tiny, secondary!60!black, align=center] at (axis cs:1.5,1.0) {underfit};
\node[font=\tiny, accent!80!black, align=center] at (axis cs:8.2,2.0) {overfit};

% --- learning curve
\nextgroupplot[title={Learning curve: error against training set size},
  xlabel={training examples $\longrightarrow$}, ylabel={error},
  xmin=0, xmax=10, ymin=0, ymax=10, xtick=\empty, ytick=\empty,
  legend style={at={(0.99,0.99)}, anchor=north east}]
\addplot[primary, line width=1.3pt, domain=0.5:9.6, samples=80] {1.1+0.9*ln(x+1)};
\addlegendentry{training}
\addplot[accent, line width=1.3pt, domain=0.5:9.6, samples=80] {9.3-2.2*ln(x+1)};
\addlegendentry{validation}
\draw[decorate, decoration={brace, amplitude=4pt, mirror}, gray!65]
     (axis cs:9.0,3.5) -- (axis cs:9.0,6.1)
     node[midway, right=4pt, font=\tiny, text=gray!55!black, align=left] {persistent gap\\= overfitting;\\more data will help};
\end{groupplot}
\end{tikzpicture}}%
{Two diagnostic plots you should produce for every model. Left: where a given complexity
sits on the trade-off. Right: whether more data would help --- a gap that is still closing
says collect more; two curves that have converged high say the model is too
simple.}{learning-curves}

\begin{definitionbox}[The Bias--Variance Trade-off]
\term{Bias} is error from wrong assumptions --- the model is too rigid to represent
the truth. \term{Variance} is error from sensitivity to the particular training
sample --- fit it again on different data and you get a very different model. Total
expected error decomposes as
\[
\text{Error} = \text{Bias}^2 + \text{Variance} + \text{irreducible noise}
\]
Simple models have high bias and low variance; complex models the reverse.
\end{definitionbox}

\rowcolors{2}{white}{lightbg}
\begin{longtable}{@{}p{3.4cm}p{5.0cm}p{5.9cm}@{}}
\toprule
\rowcolor{primary!15}
\textbf{Symptom} & \textbf{Diagnosis} & \textbf{What to do} \\
\midrule
\endhead
Train error high, validation error similarly high & Underfitting (high bias) & More complex model, better features, train longer, less regularisation \\
Train error low, validation error much higher & Overfitting (high variance) & More data, simpler model, more regularisation (\chref{features}), early stopping \\
Both low and close & Good fit & Stop. Verify once on the test set \\
Validation \emph{better} than train & Suspect a bug & Usually a leak or a mis-specified split \\
\bottomrule
\end{longtable}

\section{Baselines and the No Free Lunch Theorem}

\begin{conceptbox}[Always Build the Baseline First]
Before any model, compute what a trivial rule achieves:
\begin{tight}
  \item \emph{Classification:} always predict the majority class.
  \item \emph{Regression:} always predict the mean.
  \item \emph{Time series:} predict that tomorrow equals today.
  \item \emph{Domain rule:} whatever the humans currently do.
\end{tight}
A model that does not beat all four is not a result. This takes five minutes and
has saved more student projects than any algorithm.
\end{conceptbox}

\begin{definitionbox}[No Free Lunch]
Averaged over all possible problems, no algorithm outperforms any other. There is
no universally best model --- only models better suited to the structure of a
particular problem. This is why \chref{evaluation} exists, and why ``which
algorithm is best?'' is not a well-formed question.
\end{definitionbox}

%---------------------------------------------------------------------
\begin{fourlenses}{Setting Up a Learning Problem}
\lensLA{\emph{Paradigm:} supervised classification --- labels exist because past
outcomes are recorded. \emph{Split:} by student, and stratified, since failures are
the minority. \emph{Baseline:} ``flag anyone who missed two labs'', which is what
staff already do; beat that or the model adds nothing. \emph{Overfitting risk:}
high, because $n$ is small (one cohort) and $p$ is large (hundreds of clickstream
features).}

\lensPM{\emph{Paradigm:} supervised, but with painfully few rows --- a team may have
40 completed sprints, ever. \emph{Split:} chronological, never random: predicting
March from April is meaningless. \emph{Baseline:} last sprint's velocity.
\emph{Overfitting risk:} severe. With $n=40$, prefer a two-feature model you can
explain over a tuned gradient-boosting machine you cannot.}

\lensIOT{\emph{Paradigm:} supervised for predictive maintenance if failure records
exist; unsupervised anomaly detection if they do not --- which is the common case,
since well-maintained equipment rarely fails. \emph{Split:} by device \emph{and}
chronologically. \emph{Baseline:} fixed-interval servicing. \emph{Data:} abundant
rows but very few positive labels, so the effective sample size is the failure
count, not the row count.}

\lensSEC{\emph{Paradigm:} supervised where labelled attacks exist, unsupervised
otherwise; in practice both, side by side. \emph{Split:} strictly chronological ---
training on 2026 attacks to detect 2025 ones proves nothing. \emph{Baseline:} the
existing signature rules. \emph{Special difficulty:} the data distribution is
actively changed by an adversary, so the i.i.d.\ assumption underlying all of Part
III is false here. \chref{cybersecurity} confronts this directly.}
\end{fourlenses}

%---------------------------------------------------------------------
\begin{lab}[Split, Baseline, Diagnose]
\begin{lstlisting}[language=Python]
import numpy as np
from sklearn.model_selection import train_test_split, learning_curve
from sklearn.dummy import DummyClassifier
from sklearn.tree import DecisionTreeClassifier
import matplotlib.pyplot as plt

X, y = ...   # your features and labels

# 1. Stratified split -- keeps the rare class present everywhere.
X_tr, X_te, y_tr, y_te = train_test_split(
    X, y, test_size=0.2, stratify=y, random_state=42)

# 2. BASELINE FIRST. Never skip this.
base = DummyClassifier(strategy="most_frequent").fit(X_tr, y_tr)
print(f"baseline accuracy: {base.score(X_tr, y_tr):.3f}")

# 3. Overfitting demonstration: an unconstrained tree memorises.
deep = DecisionTreeClassifier().fit(X_tr, y_tr)
print(f"deep tree  train {deep.score(X_tr, y_tr):.3f}")   # ~1.000 -- memorised
shallow = DecisionTreeClassifier(max_depth=3).fit(X_tr, y_tr)
print(f"depth-3    train {shallow.score(X_tr, y_tr):.3f}")

# 4. Learning curve: would more data help?
sizes, tr_sc, va_sc = learning_curve(
    DecisionTreeClassifier(max_depth=5), X_tr, y_tr, cv=5,
    train_sizes=np.linspace(0.1, 1.0, 8))
plt.plot(sizes, tr_sc.mean(1), label="training")
plt.plot(sizes, va_sc.mean(1), label="validation")
plt.xlabel("training examples"); plt.ylabel("accuracy"); plt.legend(); plt.show()
\end{lstlisting}

\textbf{Read the curve:} still converging $\Rightarrow$ collect more data.
Converged with a gap $\Rightarrow$ regularise. Converged with no gap but low
$\Rightarrow$ the model is too simple or the features are inadequate.
\end{lab}

%---------------------------------------------------------------------
\begin{checkpoint}
\begin{tightnum}
  \item A model scores 0.99 on training and 0.61 on validation. Name the problem
        and give two remedies.
  \item Why can the test set be used only once, and what exactly goes wrong on the
        twentieth use?
  \item You have 400 sensor readings from each of 10 machines. Should you split by
        row or by machine? Justify.
\end{tightnum}
\end{checkpoint}

%---------------------------------------------------------------------
\begin{chaptersummary}
\begin{tight}
  \item Learning means \emph{generalising}, not memorising. Everything here
        protects that distinction.
  \item Choose the paradigm by asking whether examples of the right answer exist;
        labelling cost usually decides the project.
  \item Train fits parameters, validation chooses models, test is used once.
        Stratify for imbalance and split by entity when rows are grouped.
  \item Overfitting is low train / high validation error; underfitting is both
        high. Learning curves tell you which, and whether more data would help.
  \item Bias--variance: simple models are rigid and stable, complex ones are
        flexible and jumpy. The best model is at the balance point, not at either
        end.
  \item Build the baseline before the model. No free lunch means there is no
        universally best algorithm.
\end{tight}
\end{chaptersummary}

\begin{reviewq}
  \item \textbf{(MCQ)} High training accuracy with much lower validation accuracy
        indicates: (a)~underfitting (b)~overfitting (c)~high bias (d)~leakage only.
  \item \textbf{(MCQ)} A simple, rigid model typically has: (a)~high bias, low
        variance (b)~low bias, high variance (c)~both high (d)~both low.
  \item \textbf{(Short)} Explain the role of the validation set and why it is
        needed given that a test set already exists.
  \item \textbf{(Short)} State two sensible baselines for a regression problem and
        one for a time series.
  \item \textbf{(Short)} What does a learning curve whose two lines have converged
        at a high error tell you? What would you do next?
  \item \textbf{(Applied)} You have 40 sprints of project data and 60 candidate
        features. Explain the specific overfitting danger, and describe the model
        and validation strategy you would choose instead of a complex learner.
  \item \textbf{(Applied)} A colleague reports 99.4\% accuracy detecting network
        intrusions where 0.6\% of traffic is malicious. Explain why this number is
        uninformative, name the baseline it must be compared against, and state
        what you would ask to see instead.
\end{reviewq}
